\chapter{Parent-origin общего энергетического масштаба charged mediator}
\label{ch:v8-charged-mediator-common-energy-origin}

\section{Резонанс задаёт равенство, но не единицу энергии}

Минимальный пакет предыдущего гейта содержит общий квант
\begin{equation}
 E_*=E_C=7\gamma>0.
 \label{eq:v8-common-energy-identification}
\end{equation}
Он одновременно определяет
\begin{equation}
 g=\chi E_*,\qquad
 \tau_C=\frac{\hbar}{E_*},\qquad
 \Gamma=\chi^2\frac{E_*}{\hbar}.
 \label{eq:v8-common-energy-derived-data}
\end{equation}
Однако для всякого $\lambda>0$ преобразование
\begin{equation}
 E_*\mapsto\lambda E_*,\qquad
 t\mapsto t/\lambda
 \label{eq:v8-common-energy-scale-orbit}
\end{equation}
сохраняет безразмерное время:
\begin{equation}
 \frac{(\lambda E_*)(t/\lambda)}{\hbar}
 =\frac{E_*t}{\hbar}.
 \label{eq:v8-common-energy-time-invariant}
\end{equation}
При фиксированном $\chi$ одновременно
\begin{equation}
 g\mapsto\lambda g,\qquad
 \Gamma\mapsto\lambda\Gamma,\qquad
 \frac{\Gamma}{E_*/\hbar}=\chi^2.
 \label{eq:v8-common-energy-rate-orbit}
\end{equation}
Следовательно, resonance сокращает число координат, но не разрывает
масштабную орбиту.

\section{Локальный parent выбирает vacuum, но не gap}

Нормированный cell-projector имеет вид
\begin{equation}
 P_{\rm exc}=I_{45}-|0\rangle\langle0|.
 \label{eq:v8-common-energy-excitation-projector}
\end{equation}
Два гамильтониана
\begin{equation}
 H_1=E_*P_{\rm exc},
 \qquad H_2=2E_*P_{\rm exc}
 \label{eq:v8-common-energy-projector-witnesses}
\end{equation}
имеют одно и то же одномерное ядро,
\begin{equation}
 \ker H_1=\ker H_2=\mathbb C|0\rangle,
 \label{eq:v8-common-energy-same-vacuum}
\end{equation}
но разные спектральные щели. Поэтому state-selection и energy-selection
являются различными задачами.

\section{Исчерпание внутренних кандидатов}

Ранее проверенные величины распадаются на шесть классов:
\begin{equation}
 \mathcal C_E=\{\Lambda,R^{-1},\operatorname{gap}\mathcal L_{42},
 \operatorname{gap}P_{\rm exc},\gamma_{ud},m_{\rm obs}\}.
 \label{eq:v8-common-energy-candidates}
\end{equation}
Спектральное обрезание и радиус сохраняют собственные орбиты
перенормировки; gap генератора и projector-gap безразмерны до выбора
коэффициента; $\gamma_{ud}$ есть искомая координата, а не независимый
источник; наблюдаемая масса является внешним target-loaded данным.
Ни один кандидат не задаёт типизированной карты
\begin{equation}
 \mathfrak E_{\rm old}\longrightarrow E_*
 \quad\text{одновременно в }H_C,H_E,G_g.
 \label{eq:v8-common-energy-missing-map}
\end{equation}
Точный реестр равен
\begin{equation}
 N_{\rm candidate\ origin}=0/6.
 \label{eq:v8-common-energy-candidate-ledger}
\end{equation}
При этом масштабная орбита проверена полностью:
\begin{equation}
 N_{\rm scale\ orbit}=6/6.
 \label{eq:v8-common-energy-orbit-ledger}
\end{equation}

\section{Граница между Томами VIII и IX}

Том VIII начинался с вопроса о полном корреляционном операторе, его
локальности, марковской динамике, неподвижной алгебре и межсекторных
каналах. Все эти объекты теперь предъявлены или получили точные no-go
границы. Вывод $E_*$ требует уже не анализа данного процесса, а нового
единого parent, одновременно выбирающего четыре слота
\begin{equation}
 \mathcal D_{IX}=\{\mathfrak e_{\rm endpoint},E_*,\chi,
 \mathfrak t_{\rm transport}\}.
 \label{eq:v8-common-energy-tome9-data}
\end{equation}
Это новая центральная задача, поэтому положительная конструкция относится
к Тому IX. Текущий energy-scale gate остаётся в Томе VIII как его последний
closure-тест:
\begin{equation}
 N_{E_*\rm\ origin}=0/6,\qquad
 \text{scope}(\text{next constructive gate})=\text{Tome IX}.
 \label{eq:v8-common-energy-final-verdict}
\end{equation}

\begin{theorem}[Масштабная граница Тома VIII]
Резонансная идентификация $E_*=E_C=7\gamma$ согласует clock, mediator и
binary gap, но сохраняет непрерывную орбиту общего пересчёта энергии и
времени. Ни один из шести внутренних кандидатов не даёт типизированного
parent-origin $E_*$. Поэтому абсолютные time и rate остаются условными, а
положительное построение четырёхслотового parent переносится в Том IX.
\end{theorem}

\begin{nogocriterion}[Запрет продолжения Тома VIII подбором масштаба]
Нельзя завершать корреляционно-марковскую программу присваиванием
$E_*=m_{\rm obs}$, $E_*=\Lambda$ или $E_*=1/R$ без единого parent-морфизма.
Такой выбор не является новым следствием полного ядра и должен считаться
внешней калибровкой.
\end{nogocriterion}

\begin{protocol}\begin{enumerate}
\item common quantum: \textbf{$E_*=E_C=7\gamma$};
\item resonance выбирает значение $E_*$: \textbf{нет};
\item непрерывная scale-orbit: \textbf{есть};
\item vacuum совпадает при $H_1,H_2$: \textbf{да};
\item gap совпадает при $H_1,H_2$: \textbf{нет};
\item проверенных candidate classes: \textbf{$6$};
\item common energy parent-origin: \textbf{$0/6$};
\item absolute time/rate: \textbf{условны};
\item этот гейт относится к Тому VIII: \textbf{да, как final closure};
\item следующая конструктивная программа относится к Тому IX: \textbf{да};
\item следующий гейт: \textbf{admission четырёхслотовой программы Тома IX}.
\end{enumerate}\end{protocol}

Машинный сертификат сохранён в
\path{s2t_v8_baryon_c0_charged_mediator_common_energy_scale_parent_origin_gate_results.json}.